\begin{figure}[htbp]
	\centering
	\begin{tikzpicture}[>=Stealth,
		x={(0.85cm,0cm)}, y={(0cm,0.85cm)}, z={(0.425cm,0.255cm)},% proiezione obliqua (ax=0.5, ay=0.30) a scala 0.85
		vec/.style={->,thick},
		aux/.style={->,thin,gray!60},
		lact/.style={gray!55,thin,dashed},
		pt/.style={circle,fill,inner sep=1.1pt}]
		% ---- punti principali (coordinate: orizz, vert, profondita' [recede in alto a dx]) ----
		\coordinate (O)   at (0, 0, 0);            % O = r intersezione sigma
		\coordinate (P)   at (2.6, 0, 0.9);        % punto di applicazione, NEL piano
		\coordinate (Pq)  at (3.1, 0, 2.0);        % P + v_perp (nel piano)
		\coordinate (Pv)  at (3.1, 1.0, 2.0);      % P + v
		\coordinate (Hf)  at (1.816, 0, -0.825);   % piede del braccio d
		\coordinate (La)  at (1.55, 0, -1.41);     % linea d'azione di v_perp, estremo a
		\coordinate (Lb)  at (3.175, 0, 2.165);    % estremo b
		\coordinate (Ra)  at (0, 1.1, 0);          % O + a_r
		\coordinate (Re)  at (0, 2.6, 0);          % estremo alto di r
		% ---- piano \pi per P ortogonale a r ----
		\draw[gray!60] (-2.1, 0, -2.6) -- (4.5, 0, -2.6) -- (4.5, 0, 2.8) -- (-2.1, 0, 2.8) -- cycle;
		%\draw[gray!60] (-2.1, 0, -2.1) -- (4.5, 0, -2.1) -- (4.5, 0, 2.8)	-- (-2.1, 0, 2.8) -- cycle;
		\node[gray!75] at (-1.6, 0, -1.8) {$\pi$};
		% ---- retta r (perfora il piano in O) e versore ----
		\draw[line width=0.6pt] (0, -1.0, 0) -- (Re) node[above=1pt] {$r$};
		\coordinate (raA1) at ($(O)!2.4mm!(Re)$);
		\coordinate (raA3) at ($(O)!2.4mm!(Hf)$);
		\draw[line width=0.5pt] (raA1) -- ($(raA1)+(raA3)-(O)$) -- (raA3);
		\draw[vec] (O) -- (Ra) node[pos=0.95,anchor=east,xshift=-3pt] {$\ab_r$};
		
		% ---- verso positivo di rotazione attorno a r (mano destra, antiorario da sopra) ----
		\draw[thin, postaction={decorate},
		decoration={markings, mark=at position 1 with {\arrow{Stealth}}}]
		plot[domain=15:255,samples=90] ({0.52*cos(\x)}, 0.5, {0.52*sin(\x)});
		
		
		% ---- linea d'azione di v_perp nel piano e braccio d ----
		\draw[lact] (La) -- (Lb);
		\draw[thin] (O) -- (Hf) node[midway,anchor=south,yshift=2pt,font=\small] {$d$};
		\coordinate (raB1) at ($(Hf)!2.4mm!(Lb)$);
		\coordinate (raB3) at ($(Hf)!2.4mm!(O)$);
		\draw[line width=0.5pt] (raB1) -- ($(raB1)+(raB3)-(Hf)$) -- (raB3);
		% ---- decomposizione di v in P: v_perp nel piano, v_par lungo r ----
		\draw[aux] (Pq) -- (Pv) node[midway,anchor=west,xshift=2pt,gray!75,font=\footnotesize] {$\vb_\parallel$};
		\draw[vec] (P) -- (Pv) node[pos=0.55,anchor=south east,xshift=-1pt] {$\vb$};
		\draw[vec] (P) -- (Pq) node[pos=0.55,anchor=north west,xshift=2pt,yshift=-1pt] {$\vb_\perp$};
		% ---- punti ----
		\node[pt] at (O) {};
		\node[anchor=north east,xshift=-2pt] at (O) {$O$};
		\node[pt] at (P) {};
		\node[below right=1pt] at (P) {$P$};
		\node[pt] at (Hf) {};
		\node[anchor=north west,xshift=2pt,yshift=1pt] at (Hf) {$H$};
	\end{tikzpicture}
	\caption[Momento assiale come momento piano]{Riduzione al piano del
		momento assiale: per la doppia invarianza si sceglie il piano
		$\pi$ per $P$ ortogonale a $r$ e, come polo, l'intersezione
		$O = r \cap \pi$. Decomposto $\vb = \vb_\parallel + \vb_\perp$,
		la componente parallela non contribuisce e $\mu_r = |\vb_\perp|\,d$
		è il momento piano di $\vb_\perp$ rispetto a $O$, positivo se la
		rotazione indotta è concorde con $\ab_r$ secondo la regola della
		mano destra.}
	\label{f:momento_assiale_piano}
\end{figure}
